\documentclass[11pt]{article}

\usepackage[a4paper,margin=1in]{geometry}
\usepackage[T1]{fontenc}
\usepackage[utf8]{inputenc}
\usepackage{amsmath,amssymb,amsthm}
\usepackage{booktabs}
\usepackage{enumitem}
\usepackage{microtype}
\usepackage{xspace}
\usepackage{hyperref}
\usepackage{url}

\hypersetup{
  colorlinks=true,
  linkcolor=blue,
  citecolor=blue,
  urlcolor=blue
}

\newcommand{\Lean}{Lean~4\xspace}
\newcommand{\Mathlib}{Mathlib\xspace}
\newcommand{\DKPS}{DKPS\xspace}
\newcommand{\CMDS}{CMDS\xspace}
\newcommand{\repo}{\textsc{aiq-dkps-formalization}\xspace}
\newcommand{\file}[1]{\path{#1}}
\newcommand{\decl}[1]{\path{#1}}

\title{Formalizing Spectral Perturbation, MDS, and Alignment\
for Statistical Embeddings in Lean}

\author{TODO: Authors}
\date{Draft 3 -- June 2026}

\begin{document}
\maketitle

\begin{abstract}
We describe a \Lean formalization of a connected body of finite-dimensional
mathematics underlying statistical embedding methods: raw-stress
multidimensional-scaling consistency, classical multidimensional-scaling
perturbation, orthogonal alignment, measurable alignment events, and
probability-to-rate composition.  Data kernel perspective space (DKPS)
embeddings of black-box generative models provide the motivating application
and integration test, but the main contribution is the formalization of the
reusable mathematical infrastructure needed to make this theory kernel-checkable.

The development formalizes results that are standard enough to be used in
statistical and geometric arguments, but difficult to use in a proof assistant
without substantial interface design.  The verified statements make explicit
rank, gap, uniqueness, compactness, alignment, and measurability assumptions
that are often implicit in paper arguments.  The artifact builds against a
pinned \Lean and \Mathlib snapshot without proof holes.  We also record the
model-assisted workflow that produced the development, distinguishing the
verified mathematical artifact from provenance claims about individual language
model sessions.
\end{abstract}

\section{Introduction}

Many statistical embedding arguments combine several familiar ingredients:
empirical averages concentrate; pairwise dissimilarities inherit this
concentration; a matrix representation is constructed from those dissimilarities;
spectral perturbation controls an eigenspace or a truncated embedding; and a
Procrustes or orthogonal-alignment argument removes the non-identifiability of
coordinates.  This pattern appears in multidimensional scaling, kernel methods,
shape analysis, random graph embeddings, and response-based representations of
black-box models.  It is also the mathematical pattern behind recent work on data
kernel perspective space (DKPS) embeddings for generative models
\cite{duderstadt2023datakernels,helm2024tracking,acharyya2024consistent,acharyya2025concentration,helm2025inference,helm2026query}.

In prose, the components of this pattern are usually treated as well known.
Classical multidimensional scaling (CMDS) is described by double-centering a
squared dissimilarity matrix and taking a spectral truncation
\cite{torgerson1952mds,gower1966distance}.  Raw-stress MDS is described as an
optimization problem whose minimizers are stable under perturbations of the
input dissimilarities \cite{kruskal1964mds}.  Spectral perturbation and
Davis--Kahan-style estimates are used to compare eigenspaces
\cite{davis1970rotation,yu2015davis,chen2021spectral}.  Procrustes alignment is
used to compare configurations only up to rigid motion \cite{goodall1991procrustes}.

These facts are standard, but standard does not mean directly available in a
proof assistant in the form required by an applied theorem stack.  The theorem
prover requires explicit types, hypotheses, and measurable events.  Coordinate
representations of eigenspaces depend on choices of eigenvectors.  MDS
configurations are identifiable only up to isometry.  Minimized objectives may
have nonunique minimizers.  Probability statements cannot quantify over
nonmeasurable choices.  Downstream theorems can use an upstream result only if
its quantifiers, error notion, and measurability interface have the right shape.

This paper reports on a \Lean development that formalizes this difficult
infrastructure in the setting of a DKPS theorem stack.  DKPS is important here
because it supplies a demanding end-to-end test case: response samples are
converted to empirical dissimilarities, dissimilarities to MDS configurations,
MDS perturbation to aligned embedding error, and aligned embedding error to
query-efficiency and inference-transfer theorems.  The focus of the paper,
however, is not DKPS as a statistical methodology.  The focus is the successful
formalization of the mathematical bridges that make such a stack verifiable.

\paragraph{Contributions.}
The development makes the following contributions.
\begin{enumerate}[leftmargin=*]
  \item It formalizes a deterministic CMDS perturbation theorem converting
  matrix perturbation into aligned configuration error.  This requires formal
  spectral perturbation, rank-gap bookkeeping, overlap estimates, polar-factor
  control, orthogonal alignment, and Gram-realization lemmas.
  \item It formalizes raw-stress MDS stability through minimizer-set arguments,
  compactness, approximate-minimizer stability, and explicit uniqueness
  hypotheses for fixed-profile conclusions.
  \item It provides a proof-assistant-friendly treatment of alignment and
  measurability, replacing a noncanonical measurable choice of alignments by a
  choice-free existential event and compact-existential measurability.
  \item It composes response-level moment assumptions into high-probability
  aligned embedding-error rates, and demonstrates that the resulting interface
  supports downstream query-efficiency and inference theorems.
  \item It stages a collection of reusable lemmas in a \Mathlib-oriented
  library, including results in spectral theory, finite-dimensional geometry,
  convergence in measure, and sample-mean variance.
\end{enumerate}

The artifact was produced by a model-assisted workflow.  The paper therefore
also records provenance at the level appropriate for a formalization paper: the
mathematical result is the checked \Lean artifact; the model-assistance claims
are historical claims about how the artifact was constructed and should be
supported by repository comments, git history, and author-maintained session
records.

\section{Artifact Overview}

The repository is organized into five Lean libraries.  Four are paper-facing
application layers, and one stages reusable mathematical infrastructure.
Table~\ref{tab:libraries} summarizes their roles.

\begin{table}[t]
\centering
\begin{tabular}{@{}ll@{}}
\toprule
Library & Formalized role \\
\midrule
\file{Acharyya2024} & Raw-stress MDS and asymptotic consistency interfaces. \\
\file{Acharyya2025} & Finite-sample concentration and CMDS perturbation. \\
\file{DkpsQuench} & Query-efficiency reduction consuming aligned embedding error. \\
\file{Helm2025} & Inference-transfer reduction consuming alignment consistency. \\
\file{ForMathlib} & Reusable spectral, geometric, probability, and measurability lemmas. \\
\bottomrule
\end{tabular}
\caption{Main libraries in the development.  The DKPS layers serve as the
application stack; the formalization contribution is concentrated in the
bridges and reusable infrastructure.}
\label{tab:libraries}
\end{table}

The draft described here was validated against repository commit
\texttt{c8cf511fc4c7b7911294279b99fd76b514cbf64c}, dated 2026-06-12
10:49:10 -0400, with toolchain \texttt{leanprover/lean4:v4.31.0-rc2}.  The
inspected archive contains 51 Lean source files and 12,888 lines of Lean source.
A user-observed \texttt{lake build} at this commit completed successfully with
8,626 build jobs.  The paper directory includes a separate
\file{claims_to_files.md} index mapping claims in this draft to repository
files.

\section{Formalizing Standard Mathematics in a Nonstandard Interface}

The main technical work lies in translating familiar mathematical claims into
interfaces that can be checked and composed.  This section describes the
principal obstacles.

\subsection{Non-identifiability of configurations}

MDS and spectral embeddings do not produce canonical coordinates.  A centered
configuration is identifiable from its Gram matrix only up to an orthogonal
transformation, and eigenspaces are represented by nonunique eigenbases.  A
paper proof can often write ``after alignment'' and proceed.  A proof assistant
requires a precise statement: either an explicit isometry must be supplied, or
the event of existence of such an isometry must be made measurable.

The formalization therefore uses configuration error modulo isometry as a basic
interface.  The principal deterministic theorem in
\file{Acharyya2025/ConfigPerturbation.lean},
\decl{exists_isometry_configError_spectralConfig_le}, states that under explicit
spectral hypotheses a perturbed CMDS configuration can be aligned to the
population configuration with a quantitative error bound.  This is the theorem
that turns matrix perturbation into a statement usable by statistical arguments.

\subsection{Spectral perturbation with rank and gap hypotheses}

The CMDS bridge is not a direct application of a single library theorem.  It
requires a sequence of finite-dimensional facts.  Entrywise matrix control is
converted to operator control.  Weyl-type inequalities control eigenvalues.
Rank and eigenvalue-floor assumptions produce the spectral separation needed for
Davis--Kahan estimates.  Cross-block overlap estimates control the deviation of
empirical and population eigenspaces.  A quantitative polar-factor theorem turns
an approximate isometry into a genuine isometry.  Gram-rigidity lemmas
then transfer this information back to configurations.

These pieces are distributed across
\file{Acharyya2025/Weyl.lean}, \file{DavisKahan.lean},
\file{RankGap.lean}, \file{Overlap.lean}, \file{PolarFactor.lean},
\file{GramRigidity.lean}, \file{GramRealization.lean},
\file{OperatorBridge.lean}, and \file{ConfigPerturbation.lean}.  Several of the
same ideas are also staged in \file{ForMathlib}.  The important point is that
well-known spectral reasoning had to be decomposed into formally usable lemmas
with explicit constants, dimensions, and hypotheses.

\subsection{MDS minimizers and uniqueness}

Raw-stress MDS presents a different obstacle: the embedding is defined as a
minimizer of a nonconvex objective.  The formalized argument must establish the
right compact or coercive setting for minimizer existence, handle translation
normalization, reason about approximate minimizers, and state stability in a way
that does not pretend a nonunique minimizer is canonical.

The \file{Acharyya2024} library separates an unconditional minimizer-set
stability theorem from fixed-profile conclusions.  The fixed-profile version
requires an explicit \decl{UniquePairProfile} hypothesis.  This is not merely a
Lean technicality.  It is the formal version of a mathematical distinction that
is easy to suppress in prose: without uniqueness, the stable object is a set of
minimizers, not a chosen target configuration.  The development thereby
identifies a latent assumption and records it in the theorem interface.

\subsection{Measurability and selection}

Measurability is often where informal alignment arguments fail to translate
directly.  A deterministic theorem may prove that for each sample point there
exists an aligning isometry, but probability statements require the event being
measured to be measurable.  Selecting eigenbases or alignments by classical
choice does not automatically provide such measurability.

The development resolves a significant alignment-selection problem by changing
the event.  In \file{Acharyya2025/AlignedPipeline.lean}, the aligned estimator
error event is related to a choice-free existential predicate,
\decl{AlignExists}.  Its measurability is then proved using a compact-existential
measurability result staged in \file{ForMathlib/MeasureTheory/CompactExists.lean}.
This avoids the need for a measurable choice of an aligning isometry.

The remaining primitive, recorded in the repository as \decl{hmeas_spec}, is
smaller and more explicit.  It concerns measurability of the raw spectral
configuration map as a function of the input matrix.  Since the current
implementation uses a noncanonical eigenbasis, this is a genuine analytic seam.
The repository documents possible solutions through projector- or Gram-valued
reformulations, but leaves the primitive explicit in the relevant bridge.

\section{Principal Formalized Results}

\subsection{The deterministic CMDS perturbation theorem}

The central deterministic theorem controls CMDS configurations under matrix
perturbations.  Informally, let \(B\) be the population centered matrix and
\(\hat B\) an empirical perturbation.  Suppose \(B\) is positive semidefinite of
rank at most \(d\), its leading eigenvalues satisfy appropriate lower bounds,
and \(\hat B\) is sufficiently close to \(B\).  Then there exists an isometry of
\(\mathbb{R}^d\) aligning the empirical spectral configuration with the
population configuration, with an explicit configuration-error bound.

The formal statement is represented by declarations such as
\decl{exists_isometry_configError_spectralConfig_le} and
\decl{exists_isometry_configError_le_of_entrywise_close}.  These declarations
are the deterministic core of the concentration layer: all probabilistic
embedding-error statements eventually route through this bridge.

\subsection{Raw-stress MDS stability}

The raw-stress development includes declarations such as
\decl{rawStress_mds_stability_set}, \decl{rawStress_mds_stability},
\decl{mds_stability_inProbability_set}, and
\decl{mds_stability_inProbability_of_uniqueProfile}.  These theorems verify the
stability of MDS minimizers and approximate minimizers under perturbations of
the input dissimilarity structure.  The in-probability versions are organized
around event inclusions rather than a measurable selection of minimizers.

The result is a more robust formal account than a direct transcription of a
paper statement.  It distinguishes the unconditional minimizer-set theorem from
fixed-profile corollaries requiring uniqueness, and it records the probability
interface needed by later consistency theorems.

\subsection{Probability-to-rate composition}

The probability layer starts from response-level second-moment assumptions and
proves sample-mean concentration using Chebyshev and union-bound arguments.
These bounds imply uniform response-mean closeness, which implies closeness of
empirical dissimilarities and centered CMDS matrices.  The deterministic bridge
then yields aligned configuration error.  The rate-chain file proves that the
resulting bound tends to zero under the intended regime.

The relevant files include \file{Acharyya2024/SecondMoment.lean},
\file{Acharyya2024/Probability.lean}, \file{Acharyya2025/Bridge.lean},
\file{Acharyya2025/AlignedPipeline.lean}, and \file{Acharyya2025/RateChain.lean}.
The important formal achievement is compositionality: the downstream theorems do
not merely assume a black-box concentration event, but can be instantiated from
more primitive response-level assumptions.

\subsection{Downstream reductions as integration tests}

The \file{DkpsQuench} and \file{Helm2025} libraries are best viewed in this
paper as integration tests for the mathematical interfaces.  The query-efficiency
layer consumes a high-probability uniform embedding-error event; the Acharyya
bridge supplies this event from aligned spectral concentration, and further
variants push the input down toward response-mean and second-moment assumptions.
The Helm layer consumes an alignment-consistency condition; its bridge connects
that condition to the spectral concentration interface subject to a documented
per-sample alignment assumption.

This kind of theorem plumbing is itself significant.  A formal theorem can be
correct and still unusable if it quantifies over the wrong object or exposes the
wrong event.  The downstream reductions demonstrate that the repaired upstream
interfaces have the shape needed by later statistical theorems.

\subsection{Reusable infrastructure}

The \file{ForMathlib} layer stages general lemmas separated from the DKPS
application.  These include Courant--Fischer and Weyl lemmas, Davis--Kahan-style
bounds, Gram-matrix rigidity, quantitative near-isometry and
polar-factor results, entrywise-to-operator comparisons, compact-existential
measurability, convergence-in-measure constructors, sample-mean moment identities,
and approximate-minimizer stability.  Some are thin generalizations of local
facts; others required substantial restatement to match \Mathlib idioms, such as
generalization from real scalars to \texttt{RCLike} where appropriate.

This library makes clear that the formalization is not only about one applied
statistical construction.  It contributes a reusable substrate for future
formalizations involving spectral statistical methods and finite-dimensional
Euclidean embeddings.

\section{Model-Assisted Provenance}

The repository was developed with several language-model-assisted sessions.  The
mathematical guarantee, however, does not come from those models.  It comes from
the \Lean kernel checking the final repository against the pinned dependencies.
For that reason, provenance should be reported as an assistance statement rather
than as a substitute for verification.

The evidence currently available is heterogeneous.  Some files contain explicit
provenance comments.  For example, \file{Helm2025/Basic.lean} and
\file{DkpsQuench/Basic.lean} record that their initial scaffolds were generated
by Harmonic's Aristotle.  Many Acharyya scaffold or bridge files contain comments
such as ``Formalized by Codex 5.5 High, per user-observed model label.''  The
Acharyya READMEs record that the spectral bridge, aligned pipeline, rate chain,
and retirement of false-as-written seams were formalized by Claude Fable 5.
Several \file{ForMathlib} files and later bridge refinements record Claude Opus
4.8 as the authoring session, and the git history shows subsequent commits for
staging, rewiring, moderate generalization, and documentation after the Fable
work.

The author-supplied provenance notes add that the project began with GPT 5.2 and
Aristotle sessions, which achieved the downstream Helm2025 and Quench reductions;
that GPT/Codex 5.5 and Opus 4.8 pushed the Acharyya layers further; that Fable
completed the most difficult formalizations; and that Opus 4.8 later performed
plumbing work, staging, and moderate generalizations.  Since GPT-authored commits
were not consistently identified as GPT-authored in commit metadata, these GPT
attributions should be checked against inline comments and external session logs
before final submission.  The accompanying \file{model_provenance.md} file records
the currently recoverable repository evidence and separates it from
user-supplied provenance.

For a final venue submission, we recommend a concise assistance statement of the
following form: the artifact was produced by a human-directed, multi-model
workflow; Aristotle and GPT-family sessions produced early downstream reductions
and scaffolds; GPT/Codex 5.5 and Opus 4.8 contributed intermediate Acharyya
interfaces and subsequent plumbing/generalization; Claude Fable 5 contributed the
hardest spectral, MDS, rate, and measurability formalizations; and all
mathematical claims are based on the final build-checked Lean sources, not on the
claimed competence of any model.  This statement should be refined after the
authors audit the complete session logs.

\section{Limitations and Honest Seams}

The artifact verifies the Lean statements in the repository, not every possible
reading of the motivating papers.  Several limitations should be made explicit.

First, the theorem statements are true relative to their stated hypotheses.  Some
of those hypotheses are stronger or more explicit than the corresponding prose
in the motivating papers.  This is one of the main outputs of the formalization,
but it also means that the paper-facing claims should be reviewed by domain
experts.

Second, the raw-stress MDS layer does not formalize every distributional version
of the original asymptotic results.  In particular, repository documentation
notes that an \(L^p\)-over-model-distribution conclusion is outside the current
scope.

Third, \decl{hmeas_spec} remains a genuine primitive in the relevant bridge.
The development avoids a measurable choice of alignments, but still assumes
measurability of the raw spectral embedding map used in the current
implementation.  Discharging this primitive likely requires a projector- or
Gram-valued reformulation of the spectral embedding.

Fourth, the Helm bridge contains an explicit assumption connecting the
fixed-population spectral theorem to the per-sample population setting used by
the inference layer.  This assumption is documented as an honest seam rather
than hidden behind a theorem name.

Finally, the model-provenance record is not yet a publication-grade audit trail.
It is sufficient for a draft assistance statement, but final claims about which
model completed which component should be checked against session transcripts,
commit contents, and human review notes.

\section{Related Work}

The formalized mathematics is rooted in multidimensional scaling,
Procrustes analysis, and spectral perturbation theory
\cite{torgerson1952mds,gower1966distance,kruskal1964mds,goodall1991procrustes,davis1970rotation,yu2015davis,chen2021spectral}.
The motivating application comes from DKPS and related response-based
representations of black-box models
\cite{duderstadt2023datakernels,helm2024tracking,acharyya2024consistent,acharyya2025concentration,helm2025inference,helm2026query}.

On the formalization side, the development builds on \Lean and \Mathlib
\cite{demoura2015lean,mathlib2020}.  Its reusable layer is closest in spirit to
ordinary library development: results are extracted from an applied proof and
restated in forms suitable for reuse and possible upstreaming.  The distinctive
feature of the project is the combination of finite-dimensional spectral theory,
MDS, probability, measurability, and downstream statistical reductions in one
connected artifact.

\section{Conclusion}

This development formalizes a collection of difficult mathematical bridges
underlying statistical embedding theory.  DKPS provides the motivating theorem
stack, but the central contribution is broader: a verified treatment of CMDS
perturbation, raw-stress MDS stability, orthogonal alignment, measurable
alignment events, and response-level probability-to-rate composition in \Lean.

The artifact shows that familiar mathematical arguments may require substantial
redesign before they can be checked and composed.  Hidden uniqueness, rank, gap,
compactness, alignment, and measurability assumptions must become explicit
interfaces.  Once made explicit, these interfaces can support downstream
statistical theorems and can be factored into reusable library-style components.

As a model-assisted formalization case study, the repository is also notable for
the kind of assistance recorded in its history.  The most important assistance
was not merely local proof completion, but participation in theorem-interface
repair, bridge construction, and repository-level integration.  The final
mathematical claim remains the build-verified Lean artifact.

\begin{thebibliography}{99}

\bibitem{acharyya2024consistent}
Aranyak Acharyya, Michael W. Trosset, Carey E. Priebe, and Hayden S. Helm.
\newblock Consistent estimation of generative model representations in the data kernel perspective space.
\newblock \emph{arXiv preprint arXiv:2409.17308}, 2024.

\bibitem{acharyya2025concentration}
Aranyak Acharyya, Joshua Agterberg, Youngser Park, and Carey E. Priebe.
\newblock Concentration bounds on response-based vector embeddings of black-box generative models.
\newblock \emph{arXiv preprint arXiv:2511.08307}, 2025.

\bibitem{chen2021spectral}
Yuxin Chen, Yuejie Chi, Jianqing Fan, and Cong Ma.
\newblock Spectral methods for data science: A statistical perspective.
\newblock \emph{Foundations and Trends in Machine Learning}, 14(5):566--806, 2021.

\bibitem{davis1970rotation}
Chandler Davis and W. M. Kahan.
\newblock The rotation of eigenvectors by a perturbation. III.
\newblock \emph{SIAM Journal on Numerical Analysis}, 7(1):1--46, 1970.

\bibitem{demoura2015lean}
Leonardo de Moura, Soonho Kong, Jeremy Avigad, Floris van Doorn, and Jakob von Raumer.
\newblock The Lean theorem prover.
\newblock In \emph{Automated Deduction -- CADE-25}, LNCS 9195, pages 378--388. Springer, 2015.

\bibitem{duderstadt2023datakernels}
Brandon Duderstadt, Hayden S. Helm, and Carey E. Priebe.
\newblock Comparing foundation models using data kernels.
\newblock \emph{arXiv preprint arXiv:2305.05126}, 2023.

\bibitem{goodall1991procrustes}
Colin Goodall.
\newblock Procrustes methods in the statistical analysis of shape.
\newblock \emph{Journal of the Royal Statistical Society: Series B}, 53(2):285--321, 1991.

\bibitem{gower1966distance}
J. C. Gower.
\newblock Some distance properties of latent root and vector methods used in multivariate analysis.
\newblock \emph{Biometrika}, 53(3--4):325--338, 1966.

\bibitem{helm2024tracking}
Hayden Helm, Brandon Duderstadt, Youngser Park, and Carey E. Priebe.
\newblock Tracking the perspectives of interacting language models.
\newblock In \emph{Proceedings of the Conference on Empirical Methods in Natural Language Processing}, 2024.

\bibitem{helm2025inference}
Hayden S. Helm, Aranyak Acharyya, Youngser Park, Brandon Duderstadt, and Carey E. Priebe.
\newblock Statistical inference on black-box generative models in the data kernel perspective space.
\newblock In \emph{Findings of the Association for Computational Linguistics}, 2025.

\bibitem{helm2026query}
Hayden Helm, Ben Johnson, and Carey E. Priebe.
\newblock Query-efficient model evaluation using cached responses.
\newblock \emph{arXiv preprint arXiv:2605.07096}, 2026.

\bibitem{kruskal1964mds}
J. B. Kruskal.
\newblock Multidimensional scaling by optimizing goodness of fit to a nonmetric hypothesis.
\newblock \emph{Psychometrika}, 29:1--27, 1964.

\bibitem{mathlib2020}
The mathlib Community.
\newblock The Lean mathematical library.
\newblock In \emph{Proceedings of the 9th ACM SIGPLAN International Conference on Certified Programs and Proofs}, pages 367--381, 2020.

\bibitem{torgerson1952mds}
Warren S. Torgerson.
\newblock Multidimensional scaling: I. Theory and method.
\newblock \emph{Psychometrika}, 17(4):401--419, 1952.

\bibitem{yu2015davis}
Yi Yu, Tengyao Wang, and Richard J. Samworth.
\newblock A useful variant of the Davis--Kahan theorem for statisticians.
\newblock \emph{Biometrika}, 102(2):315--323, 2015.

\end{thebibliography}

\appendix

\section{Selected Lean Declarations}

\begin{itemize}[leftmargin=*]
  \item Raw-stress MDS stability: \decl{rawStress_mds_stability_set},
  \decl{rawStress_mds_stability}, \decl{mds_stability_inProbability_set},
  \decl{mds_stability_inProbability_of_uniqueProfile}.
  \item Acharyya2024 consistency: \decl{fixed_models_fixed_queries_consistency_of_uniqueProfile},
  \decl{fixed_models_growing_queries_consistency_of_uniqueProfile},
  \decl{growing_models_growing_queries_perStage_consistency_of_uniqueProfile}.
  \item Deterministic CMDS perturbation: \decl{exists_isometry_configError_spectralConfig_le},
  \decl{exists_isometry_configError_le_of_entrywise_close}.
  \item High-probability concentration: \decl{highProb_aligned_configError_of_entrywise_close},
  \decl{highProb_aligned_configError_of_response_mean},
  \decl{highProb_aligned_configError_endToEndRate}.
  \item Quench bridge: \decl{queryEfficient_nn_of_aligned_spectral},
  \decl{queryEfficient_nn_of_response_mean}, \decl{queryEfficient_nn_of_second_moment}.
  \item Helm bridge: \decl{alignmentConsistency_of_highProb_configError},
  \decl{alignmentConsistency_of_aligned_spectral}.
\end{itemize}

\end{document}
